% Lösungen Einheit 4 von 4 – Einheit 4 (zusammenbau v0.1)
\einheitenkopf{Einheit 4 von 4 · Einheit 4}

\erg{27}{a) Er ist der Durchschnitt der Trefferzahlen vieler Durchgänge. \quad b) höchste Säule bei 3, also $\mu = 3$; $p = 3 : 10 = 0{,}3$ \quad c) $X$: höchste Säule bei 2, $p_X = 0{,}2$; $Y$: höchste Säule bei 6, $p_Y = 0{,}6$; $Y$ hat die größere Trefferwahrscheinlichkeit, obwohl ihre höchste Säule niedriger ist \quad d) Das Maximum liegt zwischen 7 und 8, der Erwartungswert ist also $\mu = 7{,}5 = n \cdot 0{,}5$ und nicht ganzzahlig; bei geradem $n$ wäre $n \cdot 0{,}5$ ganzzahlig, also ist $n = 15$ ungerade \quad e) $\mu = 50$, $\sigma = \sqrt{100 \cdot 0{,}5 \cdot 0{,}5} = 5$; Intervall $[47{,}5; 52{,}5]$ enthält 48 bis 52; $P \approx 0{,}080 + 2 \cdot 0{,}078 + 2 \cdot 0{,}074 = 0{,}384$}
\erg{28}{a) Fehler: Tim liest $p$ an der Höhe statt an der Lage der höchsten Säule ab; sie liegt bei 4, also $\mu = 4$ und $p = 4 : 10 = 0{,}4$ \quad b) Die Säulen einer Binomialverteilung steigen bis in die Nähe des Erwartungswerts an und fallen danach ab; ist der Erwartungswert ganzzahlig, trägt er selbst die höchste Säule. \quad c) 0{,}02; 0{,}09; 0{,}23; 0{,}31; 0{,}23; 0{,}09; 0{,}02 \quad d) $\mu = 12 \cdot \frac{1}{4} = 3$, $\sigma = \sqrt{12 \cdot \frac{1}{4} \cdot \frac{3}{4}} = 1{,}5$; $\mu + 2\sigma = 6$; ab 7 richtigen Antworten}
